% Portada de J3.4: una viga en voladizo con su primer modo y la función característica con sus raíces.
\documentclass[tikz,border=0pt]{standalone}
\usepackage{resonancia}
\begin{document}
\begin{tikzpicture}
  \clip[rounded corners=6pt] (0,0) rectangle (12,3.6);
  \shade[top color=rAzul!14, bottom color=white] (0,0) rectangle (12,3.6);
  % viga en voladizo
  \fill[apoyo] (0.55,1.3) rectangle (0.9,2.9);
  \draw[trazo] (0.9,1.3) -- (0.9,2.9);
  \draw[grueso, color=rGris] (0.9,2.1) -- (5.4,2.1);
  \draw[grueso, color=rAzul, domain=0:1, samples=40, smooth, variable=\u] plot ({0.9+4.5*\u}, {2.1 - 0.95*(3*\u*\u-\u*\u*\u)/2});
  \foreach \u in {0.2,0.4,0.6,0.8,1.0} {\fill[rNaranja] ({0.9+4.5*\u},{2.1 - 0.95*(3*\u*\u-\u*\u*\u)/2}) circle (2.3pt);}
  \node[etiqueta, color=rGris] at (3.1,0.95) {\small viga empotrada–libre, 5 sensores};
  % función característica
  \draw[eje] (6.4,1.75) -- (11.7,1.75);
  \draw[eje] (6.4,0.7) -- (6.4,3.2);
  \draw[color=rVerde, line width=1.3pt, domain=0.0:8.2, samples=100, smooth, variable=\t] plot ({6.4+0.6*\t}, {1.75+0.55*(cos(\t r)+1/cosh(\t))});
  \foreach \r in {1.8751, 4.6941, 7.8548} {\fill[rNaranja] ({6.4+0.6*\r},1.75) circle (2.5pt);}
  \node[etiqueta, color=rGris] at (9.0,0.95) {\small $\cos x+1/\cosh x=0$};
  \node[etiqueta, color=rGris] at (11.2,1.45) {\small $x=\beta L$};
\end{tikzpicture}
\end{document}
